% Luc Destombe --- licence CC BY-NC-SA 4.0
% https://creativecommons.org/licenses/by-nc-sa/4.0/deed.fr
% Fichier autonome : compiler avec pdflatex, deux fois (signets, nombre de pages).
\documentclass[11pt,a4paper]{article}
\makeatletter
\newif\ifeleve
\elevefalse

\newif\iflycee@palatino
\lycee@palatinotrue

\RequirePackage[utf8]{inputenc}
\RequirePackage[T1]{fontenc}
\RequirePackage[french]{babel}
\RequirePackage{etoolbox}

\newcommand{\ShorthandsOff}{\@ifundefined{shorthandoff}{}{\shorthandoff{;:!?}}}
\newcommand{\ShorthandsOn}{\@ifundefined{shorthandon}{}{\shorthandon{;:!?}}}
\ShorthandsOff
\AtBeginDocument{\ShorthandsOn}
\AtBeginEnvironment{tikzpicture}{\ShorthandsOff}

\RequirePackage{geometry}
\RequirePackage{fancyhdr}
\RequirePackage{lastpage}
\RequirePackage{multicol}
\RequirePackage{array}
\RequirePackage{enumitem}
\RequirePackage{xcolor}
\RequirePackage{needspace}

\RequirePackage{amsmath,amssymb}
\RequirePackage{mathpazo}
\DeclareMathSizes{10.95}{10.95}{8}{5.5}
\begingroup\catcode`\_=\active \catcode`\^=\active
\gdef_#1{\sb{\scriptscriptstyle #1}}
\gdef^#1{\sp{\scriptscriptstyle\lycee@moinsepais #1}}
\endgroup
\newcommand\lycee@moins{\mathbin{%
  \setbox\z@\hbox{$\m@th\scriptscriptstyle\lycee@moinsorig$}%
  \dimen@=.55\wd\z@ \advance\dimen@-.5pt
  \hbox to.75\wd\z@{\hss$\m@th\scriptscriptstyle\vcenter{\hbox{%
    \tikz\draw[line width=.5pt,line cap=round](0,0)--(\the\dimen@,0);}}$\hss}}}
\begingroup\lccode`\~=`\- \lowercase{\endgroup\let~}\lycee@moins
\newcommand\lycee@moinsepais{\mathcode`\-="8000 }
\AtBeginDocument{\mathchardef\lycee@moinsorig=\mathcode`\-\relax}
\AtBeginDocument{\catcode`\_=12 \mathcode`\_="8000
  \catcode`\^=12 \mathcode`\^="8000 }
\newcommand\lycee@exposants{%
  \fontdimen13\textfont2=.48\fontdimen6\textfont2
  \fontdimen14\textfont2=.48\fontdimen6\textfont2
  \fontdimen15\textfont2=.48\fontdimen6\textfont2
  \fontdimen13\scriptfont2=.48\fontdimen6\scriptfont2
  \fontdimen14\scriptfont2=.48\fontdimen6\scriptfont2
  \fontdimen15\scriptfont2=.48\fontdimen6\scriptfont2}
\every@math@size\expandafter{\the\every@math@size\lycee@exposants}
\newlength{\retraitcalculs}
\setlength{\retraitcalculs}{2em}

\RequirePackage{tikz}
\usetikzlibrary{arrows.meta,calc,babel,shapes.misc}
\RequirePackage[most]{tcolorbox}
\tcbuselibrary{skins,breakable}

\definecolor{coulDef}{HTML}{1F4E79}
\definecolor{coulProp}{HTML}{7030A0}
\definecolor{coulRem}{HTML}{C55A11}
\definecolor{coulEx}{HTML}{2E7D32}
\definecolor{coulExo}{HTML}{B03A2E}
\definecolor{coulPart}{HTML}{1F4E79}
\definecolor{coulMeth}{HTML}{1B6E4F}
\definecolor{coulCourbe}{HTML}{0B5ED7}
\definecolor{coulCourbeB}{HTML}{C00000}
\definecolor{coulCourbeC}{HTML}{2E7D32}
\definecolor{coulGrille}{HTML}{8C8C8C}
\definecolor{coulRep}{HTML}{1B6E4F}
\definecolor{coulBar}{HTML}{2E7D32}

\newcommand{\R}{\mathbb{R}}
\newcommand{\Rs}{\mathbb{R}^{*}}
\renewcommand{\le}{\leqslant}
\renewcommand{\ge}{\geqslant}
\newcommand{\vect}[1]{\overrightarrow{#1}}
\newcommand{\norme}[1]{\left\lVert #1 \right\rVert}
\newcommand{\intervalle}[2]{\left[#1\,;#2\right]}
\RequirePackage{cancel}
\renewcommand{\CancelColor}{\color{coulExo}}
\newcommand{\fois}[1]{\times{\color{coulExo}#1}}
\newcommand{\barre}[1]{\cancel{#1}}

\tikzset{grille/.style={coulGrille,line width=0.4pt}}
\newcommand{\quadrillage}[4]{\draw[grille,step=1] (#1,#2) grid (#3,#4);}
\tikzset{
  croix/.style={cross out,draw,line width=0.6pt,inner sep=0pt,minimum size=4pt},
  croixdroite/.style={croix,rotate=45,line width=0.8pt},
}
\newcommand{\pt}[3]{\path (#1) node[croix]{} node[#2,font=\small,inner sep=2.5pt]{$#3$};}

\newcommand{\lycee@repere}[4]{%
  \draw[grille,step=1] (#1,#2) grid (#3,#4);
  \draw[-{Stealth[length=2mm]},thick] (#1,0)--({#3+0.4},0) node[below left=-1pt]{$x$};
  \draw[-{Stealth[length=2mm]},thick] (0,#2)--(0,{#4+0.4}) node[below left=-1pt]{$y$};
  \node[below left,font=\scriptsize,inner sep=1.5pt] at (0,0) {$0$};}
\newcommand{\lycee@gradx}[1]{\foreach \k/\l in {#1}{%
  \draw (\k,0.07)--(\k,-0.07) node[below,font=\scriptsize,inner sep=1.5pt]{$\l$};}}
\newcommand{\lycee@grady}[1]{\foreach \k/\l in {#1}{%
  \draw (0.07,\k)--(-0.07,\k) node[left,font=\scriptsize,inner sep=1.5pt]{$\l$};}}
\newcommand{\lycee@intff}[2]{\left[#1\,;#2\right]}
\newcommand{\lycee@intfo}[2]{\left[#1\,;#2\right[}
\newcommand{\lycee@intof}[2]{\left]#1\,;#2\right]}
\newcommand{\lycee@intoo}[2]{\left]#1\,;#2\right[}
\newcommand{\lycee@ens}[1]{\left\{#1\right\}}
\AtBeginDocument{%
  \providecommand{\repere}{\lycee@repere}%
  \providecommand{\gradx}{\lycee@gradx}%
  \providecommand{\grady}{\lycee@grady}%
  \providecommand{\Cf}{\mathcal{C}\sb{\scriptscriptstyle f}}%
  \providecommand{\Cg}{\mathcal{C}\sb{\scriptscriptstyle g}}%
  \providecommand{\intff}{\lycee@intff}%
  \providecommand{\intfo}{\lycee@intfo}%
  \providecommand{\intof}{\lycee@intof}%
  \providecommand{\intoo}{\lycee@intoo}%
  \providecommand{\ens}{\lycee@ens}}

\newlist{questions}{enumerate}{3}
\setlist[questions,1]{label=\arabic*\textdegree),leftmargin=*,itemsep=2pt,topsep=3pt}
\setlist[questions,2]{label=\alph*),leftmargin=*,itemsep=1pt,topsep=2pt}
\setlist[questions,3]{label=\roman*),leftmargin=*,itemsep=1pt,topsep=1pt}
\setlist[itemize]{label=\textbullet,leftmargin=*,itemsep=2pt,topsep=3pt}

\newcommand{\besoinplace}[1]{\par\penalty\@M\begingroup
  \dimen@=#1\relax\dimen@ii=\pagegoal\advance\dimen@ii-\pagetotal
  \ifdim\dimen@>\dimen@ii\ifdim\dimen@ii>\z@\vfil\fi\break\fi\endgroup}

\newcommand{\lycee@pied}{\thepage/\pageref{LastPage}}
\newcommand{\entete}[2]{%
  \pagestyle{fancy}\fancyhf{}%
  \fancyhead[L]{\footnotesize\color{coulPart}\textbf{#1}}%
  \fancyhead[R]{\footnotesize\color{coulPart}\textbf{#2}}%
  \fancyfoot[C]{\footnotesize\lycee@pied}%
  \renewcommand{\headrulewidth}{0.4pt}%
  \renewcommand{\headrule}{\hbox to\headwidth{%
    \color{coulPart!40}\leaders\hrule height \headrulewidth\hfill}}}

\RequirePackage{ccicons}
\newcommand{\lycee@licence}{{\scriptsize\color{gray}%
  \copyright\ Luc Destombe --- \ccbyncsa\ CC BY-NC-SA 4.0}}
\newif\iflycee@licence \lycee@licencetrue
\newcommand{\sanslicence}{\lycee@licencefalse}
\AtBeginDocument{\iflycee@licence\AddToHookNext{shipout/foreground}{%
  \put(0,0){\rlap{\kern\dimexpr1in+\hoffset+\oddsidemargin+\textwidth\relax
    \llap{\raisebox{-\dimexpr1in+\voffset+\topmargin+\headheight+\headsep
      +\textheight+\footskip\relax}{\lycee@licence}}}}}\fi}

\newcommand{\formulesserrees}{%
  \setlength{\abovedisplayskip}{3pt}\setlength{\belowdisplayskip}{3pt}%
  \setlength{\abovedisplayshortskip}{2pt}\setlength{\belowdisplayshortskip}{3pt}}

\setlength{\parindent}{0pt}

\RequirePackage[implicit=false,hidelinks,pdfencoding=auto,psdextra,bookmarksopen,
  bookmarksopenlevel=1,pdfcreator={},pdfproducer={}]{hyperref}
\RequirePackage{bookmark}
\hypersetup{pdfauthor={Luc Destombe},pdfkeywords={CC BY-NC-SA 4.0}}
\newcounter{lycee@signet}
\newcommand{\signet}[3][]{\stepcounter{lycee@signet}%
  \edef\lycee@dest{\if\relax\detokenize{#1}\relax
    signet.\arabic{lycee@signet}\else#1\fi}%
  \hypertarget{\lycee@dest}{}\bookmark[level=#2,dest=\lycee@dest]{#3}}
\pdfstringdefDisableCommands{%
  \def\R{R}\def\vect#1{#1}\def\Cf{Cf}\def\Cg{Cg}\def\fois#1{×#1}%
  \def\barre#1{#1}\def\intervalle#1#2{[#1;#2]}\def\intff#1#2{[#1;#2]}%
  \def\textsuperscript#1{#1}\def\dfrac#1#2{#1/#2}\def\frac#1#2{#1/#2}%
  \def\vec#1{#1⃗}}%
\geometry{top=2cm,bottom=1cm,left=1cm,right=1cm,headheight=14pt,footskip=0.6cm}

\RequirePackage{pgfplots}
\pgfplotsset{compat=1.18}
\RequirePackage{tkz-tab}

\newcounter{partiecours}
\renewcommand{\thepartiecours}{\Roman{partiecours}}
\newcommand{\partiecours}[1]{%
  \refstepcounter{partiecours}%
  \Needspace*{8\baselineskip}%
  \bigskip
  \noindent\begin{tcolorbox}[enhanced,colback=coulPart,colframe=coulPart,
      boxrule=0pt,arc=2pt,left=8pt,right=8pt,top=4pt,bottom=4pt,
      before skip=6pt,after skip=10pt]
    \color{white}\bfseries\large \thepartiecours{} --- #1\signet{1}{\thepartiecours{} --- #1}
  \end{tcolorbox}%
  \addcontentsline{toc}{section}{\thepartiecours{} --- #1}%
}

\newtcolorbox{boitecours}[3][]{%
  enhanced,breakable,
  colback=white,colframe=#2,coltitle=white,
  fonttitle=\bfseries,
  attach boxed title to top left={xshift=8pt,yshift=-\tcboxedtitleheight/2},
  boxed title style={colback=#2,arc=2pt,boxrule=0pt},
  boxrule=0.9pt,arc=2pt,
  left=8pt,right=8pt,top=8pt,bottom=6pt,
  title={#3},#1}

\newcounter{methode}
\newenvironment{definitionbox}[1][Définition]%
  {\begin{boitecours}[phantom={\signet{2}{#1}}]{coulDef}{#1}}{\end{boitecours}}
\newenvironment{proprietebox}[1][Propriété]%
  {\begin{boitecours}[phantom={\signet{2}{#1}}]{coulProp}{#1}}{\end{boitecours}}
\newenvironment{methodebox}[1][Méthode]{\stepcounter{methode}%
  \begin{boitecours}[phantom={\signet[methode-\arabic{methode}]{2}{#1}}]{coulMeth}{#1}}%
  {\end{boitecours}}

\newcommand{\etape}[1]{\par\smallskip\textbf{\color{coulMeth}#1}\par\nobreak\smallskip}
\newcommand{\enonce}[1]{\emph{#1}\par\smallskip}
\newcommand{\reponse}[1]{\par\smallskip
  {\footnotesize\itshape\color{black!55}Réponses : #1}}
\newcommand{\demo}[1]{\textbf{\color{coulMeth}Démonstration.} #1}
\newcommand{\afaire}[1]{%
  \par\smallskip
  \noindent{\small\color{coulExo}$\blacktriangleright$\ \textbf{À faire :}
    fiche d'exercices du chapitre, #1.}\par\smallskip}

\newenvironment{calculs}{\setlength{\jot}{5pt}\@fleqntrue\@mathmargin\retraitcalculs
  \setlength{\abovedisplayskip}{4pt}\setlength{\belowdisplayskip}{4pt}%
  \setlength{\abovedisplayshortskip}{4pt}\setlength{\belowdisplayshortskip}{4pt}%
  \csname alignat*\endcsname{2}}%
  {\csname endalignat*\endcsname}
\newcommand{\rem}[1]{\text{\small\itshape\color{black!60}#1}}
\newcommand{\explique}[2]{\par\smallskip\noindent
  \begin{minipage}[t]{0.57\linewidth}\raggedright #1\end{minipage}\hfill
  \begin{minipage}[t]{0.40\linewidth}\raggedright\leavevmode
    {\small\itshape\color{black!60}#2}\end{minipage}\par}
\newcommand{\cas}[1]{\par\smallskip\noindent\textbf{\color{coulExo}#1}\nobreak}
\newcommand{\calc}[1]{\par\smallskip\textbf{\color{coulExo}#1}\par\nobreak}

\newtcolorbox{remarquebox}[1][Remarque]{%
  enhanced,breakable,colback=white,colframe=coulRem,
  boxrule=0pt,leftrule=2.5pt,arc=0pt,outer arc=0pt,
  left=8pt,right=6pt,top=5pt,bottom=5pt,
  before upper={\textbf{\color{coulRem}#1}\par\smallskip}}

\newtcolorbox{exemplebox}[1][Exemple]{%
  enhanced,breakable,colback=white,colframe=coulEx,
  boxrule=0pt,leftrule=2.5pt,arc=0pt,outer arc=0pt,
  left=8pt,right=6pt,top=5pt,bottom=5pt,
  before upper={\textbf{\color{coulEx}#1}\par\smallskip}}

\pgfplotsset{
  repere/.style={
    axis lines=middle,
    axis line style={-{Stealth[length=2.4mm]},thick},
    every axis x label/.style={at={(ticklabel* cs:1.0)},anchor=west},
    every axis y label/.style={at={(ticklabel* cs:1.0)},anchor=south},
    label style={font=\footnotesize},
    tick label style={font=\scriptsize},
    grid=both,
    major grid style={grille},
    minor tick num=0,
    tick align=inside,
    clip mode=individual,
  },
}
\tikzset{
  courbe/.style={coulCourbe,line width=1.1pt,smooth},
  courbeB/.style={coulCourbeB,line width=1.1pt,smooth},
  courbeC/.style={coulCourbeC,line width=1.1pt,smooth},
}

\newlist{etapes}{enumerate}{1}
\setlist[etapes]{label=\textbf{\arabic*.},leftmargin=*,itemsep=1pt,topsep=2pt}

\newlist{expressions}{enumerate}{1}
\setlist[expressions]{label={\Alph*\ $=$},leftmargin=*,itemsep=3pt,topsep=3pt}

\newcounter{exo}
\renewcommand{\theexo}{\ifnum\value{exo}<10 0\fi\arabic{exo}}
\newtcolorbox{exercice}[1][]{%
  enhanced,breakable,
  colback=white,colframe=coulExo,
  boxrule=0.9pt,arc=2pt,
  left=8pt,right=8pt,top=8pt,bottom=6pt,
  coltitle=white,fonttitle=\bfseries,
  attach boxed title to top left={xshift=8pt,yshift=-\tcboxedtitleheight/2},
  boxed title style={colback=coulExo,arc=2pt,boxrule=0pt},
  phantom={\signet{2}{Exercice \arabic{exo}}},
  title={Exercice \theexo\ifx\relax#1\relax\else{} --- #1\fi}}
\newenvironment{exo}[1][\relax]{\refstepcounter{exo}\begin{exercice}[#1]}{\end{exercice}}

  \newenvironment{correction}{\par}{\par}
  \newcommand{\autableau}{}
  \newcommand{\etapeexemples}{\etape{Exemples corrigés}}
  \newcommand{\etapetableau}[1]{\etape{#1}}
  \newenvironment{exempletableau}[1][Exemple]%
    {\begin{exemplebox}[#1]}{\end{exemplebox}}

\newcommand{\coursEntete}{}
\newcommand{\coursNiveau}{}
\newcommand{\coursTitre}{}
\newcommand{\coursSurTitre}{}
\newcommand{\coursSousTitre}{}

\newcommand{\enteteCours}[1]{\renewcommand{\coursEntete}{#1}}
\newcommand{\chapitrecours}[2]{\enteteCours{Chapitre #1 --- #2}}
\newcommand{\niveaucours}[1]{\renewcommand{\coursNiveau}{#1}}
\newcommand{\titrecours}[3][]{%
  \renewcommand{\coursTitre}{#2}%
  \renewcommand{\coursSurTitre}{#1}%
  \renewcommand{\coursSousTitre}{#3}}

\pagestyle{fancy}
\fancyhf{}
\fancyhead[L]{\footnotesize\color{coulPart}\textbf{\coursEntete}}
\fancyhead[R]{\footnotesize\color{coulPart}\textbf{\coursNiveau}}
\fancyfoot[C]{\footnotesize\thepage}
\renewcommand{\headrulewidth}{0.4pt}
\renewcommand{\headrule}{\hbox to\headwidth{\color{coulPart!40}\leaders\hrule height \headrulewidth\hfill}}

\setlength{\parskip}{2pt}

\AtBeginDocument{%
  \begin{center}
    \begin{tcolorbox}[enhanced,width=0.72\linewidth,colback=white,colframe=coulPart,
        boxrule=1.6pt,arc=3pt,halign=center,top=10pt,bottom=10pt]
      \ifx\coursSurTitre\empty
        {\Huge\bfseries\color{coulPart} \coursTitre}\\[4pt]
      \else
        {\Huge\bfseries\color{coulPart} \coursTitre}\\[3pt]
        {\Large\bfseries\color{coulPart} \coursSurTitre}\\[5pt]
      \fi
      {\normalsize \coursSousTitre}%
      
    \end{tcolorbox}
  \end{center}%
}
\makeatother

\chapitrecours{3}{Polynômes du second degré}
\niveaucours{1\textsuperscript{re} STI2D}
\titrecours{POLYNÔMES DU SECOND DEGRÉ}{Cours --- Première STI2D}

\begin{document}

\begin{remarquebox}[Comment utiliser ce cours]
Chaque partie commence par ce qu'il faut \textbf{savoir} (définitions, théorèmes), puis
vient ce qu'il faut \textbf{savoir faire} : une méthode, avec des
\textbf{exemples corrigés} faits en classe, puis des exercices
\og\textbf{À vous de jouer}\fg{} à chercher seul.

\end{remarquebox}

\partiecours{Fonction polynôme du second degré}

\begin{definitionbox}[Définition 1 --- Fonction polynôme du second degré]
Une \textbf{fonction polynôme de degré 2}, ou \textbf{trinôme}, est une fonction $f$
définie sur $\R$ par
\[
  f(x)=ax^{2}+bx+c \qquad\text{avec } a\ne 0 .
\]
Les nombres $a$, $b$ et $c$ sont ses \textbf{coefficients}. Si $a=0$, il reste $bx+c$ :
c'est une fonction affine.
\end{definitionbox}

\begin{exemplebox}[Exemples et contre-exemple]
$f(x)=2x^{2}-5x+1$ et $g(x)=7-3x^{2}$ sont des trinômes. En revanche, $h(x)=4x+1$ n'en
est pas un : il n'y a pas de $x^{2}$.
\end{exemplebox}

\begin{methodebox}[Méthode 1 --- Identifier les coefficients $a$, $b$ et $c$]
\etapeexemples
Donner, quand c'est possible, les coefficients $a$, $b$ et $c$ :
\begin{questions}[label=\alph*)]
  \item $P(x)=4x^{2}-3x+7$ \; ; \; $R(x)=-x^{2}+\dfrac{3}{4}x$ \; ; \;
        $T(x)=\dfrac{2}{3}x^{2}-5$
  \item $U(x)=(x+3)^{2}-2(x-1)$
  \item $V(x)=(x+2)^{2}-x^{2}$
\end{questions}
\begin{correction}
\cas{a)}
\explique{$P$ : $a=4$, $b=-3$, $c=7$.}{on lit chaque coefficient \textbf{avec son signe}}
\explique{$R$ : $a=-1$, $b=\dfrac{3}{4}$, $c=0$.}{$-x^{2}=-1\times x^{2}$ ; pas de
nombre seul : $c=0$}
\explique{$T$ : $a=\dfrac{2}{3}$, $b=0$, $c=-5$.}{pas de terme en $x$ : $b=0$}
\cas{b)}
\begin{calculs}
U(x) &= (x+3)^{2}-2(x-1) &\qquad& \rem{des parenthèses : on développe d'abord}\\
U(x) &= x^{2}+6x+9-2x+2 &\qquad& \rem{identité remarquable, puis distributivité}\\
U(x) &= x^{2}+4x+11 &\qquad& \rem{on réduit et on range}
\end{calculs}
\explique{$a=1$, $b=4$, $c=11$.}{}
\cas{c)}
\begin{calculs}
V(x) &= x^{2}+4x+4-x^{2} &\qquad& \rem{on développe}\\
V(x) &= 4x+4 &\qquad& \rem{les $x^{2}$ s'éliminent}
\end{calculs}
\explique{$V$ n'est pas un trinôme : c'est une fonction affine.}{$a=0$}
\end{correction}

\etape{À vous de jouer}
Donner, quand c'est possible, les coefficients $a$, $b$ et $c$ :
\[
  f(x)=3-2x^{2} \qquad g(x)=(3x-1)(x+2) \qquad h(x)=3(x-1)^{2}-3x^{2}
\]
\reponse{$f$ : $a=-2$, $b=0$, $c=3$ ;\; $g(x)=3x^{2}+5x-2$ ;\; $h(x)=-6x+3$ : pas un
trinôme.}
\end{methodebox}

\afaire{exercices 1 à 4}

\partiecours{Forme factorisée et racines}

\begin{definitionbox}[Définition 2 --- Forme factorisée, racines]
Une fonction $f(x)=a(x-x_{1})(x-x_{2})$, avec $a\ne 0$, est un trinôme écrit sous
\textbf{forme factorisée} : par exemple $3(x-1)(x+1)=3x^{2}-3$.

\smallskip
Les \textbf{racines} d'un polynôme $f$ sont les solutions de $f(x)=0$ : ce sont les
abscisses des points où sa courbe coupe l'axe des abscisses.
\end{definitionbox}

\begin{proprietebox}[Propriété 1 --- Racines et forme factorisée]
Les racines de $f(x)=a(x-x_{1})(x-x_{2})$ sont $x_{1}$ et $x_{2}$ (éventuellement
égales) : elles se lisent sans calcul, en changeant le signe. Dans $(x-4)$, la racine
est $4$ ; dans $(x+2)$, elle est $-2$.
\end{proprietebox}

\demo{Un produit est nul si et seulement si l'un de ses facteurs est nul ; comme
$a\ne 0$, $a(x-x_{1})(x-x_{2})=0$ équivaut à $x=x_{1}$ ou $x=x_{2}$.}

\begin{remarquebox}
Tous les trinômes ne se factorisent pas : $x^{2}+4$ n'a pas de racine, car
$x^{2}+4\ge 4$ pour tout $x$. Le \textbf{discriminant} (partie III) dira lesquels se
factorisent.
\end{remarquebox}

\begin{methodebox}[Méthode 2 --- Développer une forme factorisée, lire les racines]
\etapeexemples
\begin{questions}[label=\alph*)]
  \item Développer $f(x)=-2(x-3)(x+1)$ et donner ses racines.
  \item Même question pour $g(x)=3(x+2)^{2}$.
  \item Résoudre $-2(x-3)(x+1)=0$, puis $3(x+2)^{2}=0$.
\end{questions}
\begin{correction}
\cas{a)}
\begin{calculs}
f(x) &= -2\left(x^{2}+x-3x-3\right) &\qquad& \rem{on développe d'abord les deux parenthèses}\\
f(x) &= -2\left(x^{2}-2x-3\right) &&\\
f(x) &= -2x^{2}+4x+6 &\qquad& \rem{puis on multiplie par $a=-2$}
\end{calculs}
\explique{Racines : $3$ et $-1$.}{on annule chaque parenthèse}
\cas{b)}
\begin{calculs}
g(x) &= 3\left(x^{2}+4x+4\right) &\qquad& \rem{identité remarquable}\\
g(x) &= 3x^{2}+12x+12 &&
\end{calculs}
\explique{Une seule racine, $-2$ : on dit qu'elle est \textbf{double}.}{$g(x)=3(x+2)(x+2)$}
\cas{c)}
\explique{$S=\ens{-1\,;3}$ \quad et \quad $S=\ens{-2}$.}{produit nul ; le facteur $-2$
(ou $3$) ne s'annule jamais}
\end{correction}

\etape{À vous de jouer}
Développer, puis donner les racines :
\[
  f(x)=(x-4)(x+1) \qquad g(x)=-3(x+2)(x-5) \qquad h(x)=2(x-3)^{2}
\]
\reponse{$f(x)=x^{2}-3x-4$, racines $4$ et $-1$ ;\; $g(x)=-3x^{2}+9x+30$, racines $-2$ et
$5$ ;\; $h(x)=2x^{2}-12x+18$, racine double $3$.}
\end{methodebox}

\afaire{exercices 5 à 8}

\partiecours{Le discriminant}

\begin{definitionbox}[Définition 3 --- Discriminant]
Le \textbf{discriminant} du trinôme $ax^{2}+bx+c$ est le nombre
\[
  \boxed{\Delta=b^{2}-4ac}
\]
(lire « delta »). Son signe va dire si le trinôme a deux racines, une seule ou aucune.
\end{definitionbox}

\begin{methodebox}[Méthode 3 --- Calculer un discriminant]
\etapeexemples
Pour chaque trinôme $P$, $Q$, $R$, $S$, calculer le discriminant et donner son signe :
\[
  P(x)=2x^{2}+3x-2 \qquad Q(x)=x^{2}-6x+9 \qquad R(x)=-3x^{2}+2x-4 \qquad S(x)=9x^{2}-4
\]
\begin{correction}
\cas{Résolution}\par\nopagebreak
\begin{calculs}
\Delta_{P} &= 3^{2}-4\times 2\times(-2) &\qquad& \rem{$a=2$, $b=3$, $c=-2$ ; un nombre négatif entre parenthèses}\\
\Delta_{P} &= 25 &\qquad& \rem{$9+16$ : $\Delta_{P}>0$}\\[3pt]
\Delta_{Q} &= (-6)^{2}-4\times 1\times 9 &\qquad& \rem{$b^{2}$ est toujours positif : $(-6)^{2}=36$}\\
\Delta_{Q} &= 0 &&\\[3pt]
\Delta_{R} &= 2^{2}-4\times(-3)\times(-4) &\qquad& \rem{$-4ac$ : attention aux deux signes $-$}\\
\Delta_{R} &= -44 &\qquad& \rem{$4-48$ : $\Delta_{R}<0$}\\[3pt]
\Delta_{S} &= 0^{2}-4\times 9\times(-4) &\qquad& \rem{pas de terme en $x$ : $b=0$}\\
\Delta_{S} &= 144 &\qquad& \rem{$\Delta_{S}>0$}
\end{calculs}
\end{correction}

\etape{À vous de jouer}
Calculer le discriminant et donner son signe :
\[
  f(x)=4x^{2}-4x+1 \qquad g(x)=3x^{2}+2x+1 \qquad h(x)=-x^{2}+2x+8
\]
\reponse{$\Delta_{f}=16-16=0$ ;\; $\Delta_{g}=4-12=-8<0$ ;\; $\Delta_{h}=4+32=36>0$.}
\end{methodebox}

\afaire{exercices 9 à 11}

\partiecours{Racines du trinôme et équations du second degré}

\begin{proprietebox}[Théorème 1 --- Les trois cas]
Pour l'équation $ax^{2}+bx+c=0$ (avec $a\ne 0$) et $\Delta=b^{2}-4ac$ :
\begin{itemize}
  \item si $\Delta>0$ : deux solutions, $x_{1}=\dfrac{-b-\sqrt{\Delta}}{2a}$ et
        $x_{2}=\dfrac{-b+\sqrt{\Delta}}{2a}$ ;
  \item si $\Delta=0$ : une seule solution, la racine double $x_{0}=\dfrac{-b}{2a}$ ;
  \item si $\Delta<0$ : aucune solution ($\sqrt{\Delta}$ n'existe pas).
\end{itemize}
\end{proprietebox}

\begin{methodebox}[Méthode 4 --- Résoudre une équation du second degré]
\etapeexemples
Résoudre dans $\R$ :
\begin{questions}[label=\alph*)]
  \item $2x^{2}+5x-3=0$
  \item $x^{2}+6x+9=0$
  \item $2x^{2}-3x+4=0$
  \item $3x^{2}+x-5=2x^{2}+7$
  \item $x^{2}-9=0$ \; ; \; $2x^{2}+10x=0$ \; ; \; $x^{2}=5$
\end{questions}
\begin{correction}
\cas{a)}
\begin{calculs}
\Delta &= 5^{2}-4\times 2\times(-3) &\qquad& \rem{$a=2$, $b=5$, $c=-3$}\\
\Delta &= 49 &\qquad& \rem{$\Delta>0$ : deux solutions ; $\sqrt{49}=7$}\\
x_{1} &= \frac{-5-7}{4}=-3 &\qquad& \rem{$2a=4$}\\
x_{2} &= \frac{-5+7}{4}=\frac{1}{2} &&
\end{calculs}
\explique{$S=\ens{-3\,;\tfrac{1}{2}}$.}{}
\cas{b)}
\begin{calculs}
\Delta &= 6^{2}-4\times 1\times 9 &&\\
\Delta &= 0 &\qquad& \rem{une seule solution}\\
x_{0} &= \frac{-6}{2}=-3 &\qquad& \rem{$x_{0}=\frac{-b}{2a}$}
\end{calculs}
\explique{$S=\ens{-3}$.}{}
\cas{c)}
\begin{calculs}
\Delta &= (-3)^{2}-4\times 2\times 4 &&\\
\Delta &= -23 &\qquad& \rem{$9-32$ : $\Delta<0$}
\end{calculs}
\explique{$S=\varnothing$.}{aucune solution}
\cas{d)}
\begin{calculs}
3x^{2}+x-5 &= 2x^{2}+7 &\qquad& \rem{on recopie}\\
x^{2}+x-12 &= 0 &\qquad& \rem{tout dans le membre de gauche}
\end{calculs}
\begin{calculs}
\Delta &= 1^{2}-4\times 1\times(-12) &&\\
\Delta &= 49 &&\\
x_{1} &= \frac{-1-7}{2}=-4 &&\\
x_{2} &= \frac{-1+7}{2}=3 &&
\end{calculs}
\explique{$S=\ens{-4\,;3}$.}{}
\cas{e)}
\explique{Ces trois équations se résolvent sans $\Delta$.}{plus rapide et plus sûr
que le discriminant}
\begin{calculs}
x^{2}-9 &= 0 &\qquad& \rem{on recopie}\\
(x-3)(x+3) &= 0 &\qquad& \rem{différence de deux carrés}\\
x-3 &= 0 \quad\text{ou}\quad x+3=0 &\qquad& \rem{produit nul}
\end{calculs}
\explique{$S=\ens{-3\,;3}$.}{}
\begin{calculs}
2x^{2}+10x &= 0 &\qquad& \rem{on recopie}\\
2x(x+5) &= 0 &\qquad& \rem{pas de nombre seul : facteur commun $2x$}\\
2x &= 0 \quad\text{ou}\quad x+5=0 &\qquad& \rem{produit nul}
\end{calculs}
\explique{$S=\ens{-5\,;0}$.}{}
\begin{calculs}
x^{2} &= 5 &\qquad& \rem{pas de terme en $x$ : $x^{2}$ est déjà isolé}\\
x &= -\sqrt{5} \quad\text{ou}\quad x=\sqrt{5} &\qquad& \rem{$5>0$ : deux solutions}
\end{calculs}
\explique{$S=\ens{-\sqrt{5}\,;\sqrt{5}}$.}{}
\end{correction}

\etape{À vous de jouer}
Résoudre dans $\R$ :
\[
  x^{2}-7x+10=0 \qquad 3x^{2}-2x+2=0 \qquad 9x^{2}+6x+1=0 \qquad x^{2}-6x=0
\]
\reponse{$\Delta=9$ : $S=\{2\,;5\}$ ;\; $\Delta=-20$ : $S=\varnothing$ ;\; $\Delta=0$ :
$S=\left\{-\frac{1}{3}\right\}$ ;\; $x(x-6)=0$ : $S=\{0\,;6\}$.}
\end{methodebox}

\afaire{exercices 12 à 16}

\partiecours{Factoriser un trinôme}

\begin{proprietebox}[Théorème 2 --- Factorisation]
Pour le trinôme $P(x)=ax^{2}+bx+c$ et $\Delta=b^{2}-4ac$ :
\begin{itemize}
  \item si $\Delta>0$ : $P(x)=a(x-x_{1})(x-x_{2})$ ;
  \item si $\Delta=0$ : $P(x)=a(x-x_{0})^{2}$ ;
  \item si $\Delta<0$ : $P$ ne se factorise pas.
\end{itemize}
On n'oublie pas le coefficient $a$ devant les parenthèses.
\end{proprietebox}

\begin{methodebox}[Méthode 5 --- Factoriser un trinôme]
\etapeexemples
Factoriser, quand c'est possible :
\begin{questions}[label=\alph*)]
  \item $f(x)=3x^{2}-3x-6$
  \item $g(x)=-2x^{2}-7x-3$
  \item $h(x)=x^{2}+8x+16$ \; et \; $k(x)=x^{2}-2x+5$
\end{questions}
\begin{correction}
\cas{a)}
\begin{calculs}
\Delta &= (-3)^{2}-4\times 3\times(-6) &&\\
\Delta &= 81 &\qquad& \rem{$\sqrt{81}=9$}\\
x_{1} &= \frac{3-9}{6}=-1 &&\\
x_{2} &= \frac{3+9}{6}=2 &&
\end{calculs}
\explique{$f(x)=3(x+1)(x-2)$.}{contrôle : $3\times 1\times(-2)=-6$, c'est bien $c$}
\cas{b)}
\begin{calculs}
\Delta &= (-7)^{2}-4\times(-2)\times(-3) &&\\
\Delta &= 25 &&\\
x_{1} &= \frac{7-5}{-4}=-\frac{1}{2} &\qquad& \rem{$-b=7$ et $2a=-4$}\\
x_{2} &= \frac{7+5}{-4}=-3 &&
\end{calculs}
\explique{$g(x)=-2\left(x+\tfrac{1}{2}\right)(x+3)$.}{contrôle :
$-2\times\frac{1}{2}\times 3=-3$}
\cas{c)}
\explique{$h$ : $\Delta=64-64=0$, $x_{0}=-4$ et $h(x)=(x+4)^{2}$.}{$\Delta=0$ : un
carré}
\explique{$k$ : $\Delta=4-20=-16<0$ : $k$ ne se factorise pas.}{}
\end{correction}

\etape{À vous de jouer}
Factoriser, quand c'est possible :
\[
  f(x)=x^{2}+x-12 \qquad g(x)=-3x^{2}+12x-12 \qquad h(x)=2x^{2}+x+1
\]
\reponse{$\Delta_{f}=49$ : $f(x)=(x+4)(x-3)$ ;\; $\Delta_{g}=0$ : $g(x)=-3(x-2)^{2}$ ;\;
$\Delta_{h}=-7<0$ : pas de factorisation.}
\end{methodebox}

\begin{remarquebox}[Racine évidente]
Parfois une racine se voit : dans $f(x)=3x^{2}+2x-5$, $3+2-5=0$, donc $f(1)=0$. On écrit
$f(x)=3(x-1)(x-x_{2})$ ; en $x=0$ : $-5=3\times(-1)\times(-x_{2})$, donc
$x_{2}=-\frac{5}{3}$ et $f(x)=3(x-1)\left(x+\frac{5}{3}\right)$, sans calculer $\Delta$.
\end{remarquebox}

\afaire{exercices 17 à 20}

\partiecours{Représentation graphique des fonctions du second degré}

\begin{proprietebox}[Propriété 2 --- La parabole et son sens]
La courbe d'un trinôme est une \textbf{parabole}. Son sens dépend du signe de $a$ :
si $a>0$, elle est tournée vers le haut (un « creux ») ; si $a<0$, vers le bas (une
« bosse »).
\end{proprietebox}

\begin{center}
\begin{minipage}[t]{0.47\linewidth}\centering
{\small $a>0$ : $f(x)=x^{2}-4$}\par\smallskip
\begin{tikzpicture}[x=0.42cm,y=0.42cm]
  \repere{-4}{-5}{4}{5}
  \gradx{-3,-2,-1,1,2,3} \grady{-4,-2,2,4}
  \draw[coulCourbe,line width=1.1pt,domain=-3:3,samples=60] plot (\x,{\x*\x-4});
\end{tikzpicture}
\end{minipage}\hfill
\begin{minipage}[t]{0.47\linewidth}\centering
{\small $a<0$ : $g(x)=-x^{2}+2x+3$}\par\smallskip
\begin{tikzpicture}[x=0.42cm,y=0.42cm]
  \repere{-3}{-5}{5}{5}
  \gradx{-2,-1,1,2,3,4} \grady{-4,-2,2,4}
  \draw[coulCourbe,line width=1.1pt,domain=-2:4,samples=60] plot (\x,{-\x*\x+2*\x+3});
\end{tikzpicture}
\end{minipage}
\end{center}

\begin{proprietebox}[Propriété 3 --- Racines et axe des abscisses]
Les racines sont les abscisses des points où la parabole coupe l'axe des abscisses :
deux points si $\Delta>0$ ; un seul si $\Delta=0$ (la parabole touche l'axe) ; aucun si
$\Delta<0$.
\end{proprietebox}

\begin{center}
\begin{minipage}[t]{0.31\linewidth}\centering
{\small $\Delta>0$ : $x^{2}+2x-3$}\par\smallskip
\begin{tikzpicture}[x=0.45cm,y=0.45cm]
  \repere{-5}{-5}{3}{5}
  \gradx{-4,-3,-2,-1,1,2} \grady{-4,-2,2,4}
  \draw[coulCourbe,line width=1.1pt,domain=-4:2,samples=60] plot (\x,{\x*\x+2*\x-3});
  \path (-3,0) node[croix]{} (1,0) node[croix]{};
  \node[red,above right,font=\small,inner sep=1pt] at (-3,0) {$x_{1}$};
  \node[red,above left,font=\small,inner sep=1pt] at (1,0) {$x_{2}$};
\end{tikzpicture}
\end{minipage}\hfill
\begin{minipage}[t]{0.31\linewidth}\centering
{\small $\Delta=0$ : $x^{2}+2x+1$}\par\smallskip
\begin{tikzpicture}[x=0.45cm,y=0.45cm]
  \repere{-5}{-5}{3}{5}
  \gradx{-4,-3,-2,1,2} \grady{-4,-2,2,4}
  \draw[coulCourbe,line width=1.1pt,domain=-3.2:1.2,samples=60] plot (\x,{\x*\x+2*\x+1});
  \path (-1,0) node[croix]{};
  \node[red,below,font=\small,fill=white,inner sep=1pt] at (-1,-0.2) {$x_{0}$};
\end{tikzpicture}
\end{minipage}\hfill
\begin{minipage}[t]{0.31\linewidth}\centering
{\small $\Delta<0$ : $x^{2}+2x+3$, aucune racine}\par\smallskip
\begin{tikzpicture}[x=0.45cm,y=0.45cm]
  \repere{-5}{-5}{3}{5}
  \gradx{-4,-3,-2,-1,1,2} \grady{-4,-2,2,4}
  \draw[coulCourbe,line width=1.1pt,domain=-2.7:0.7,samples=60] plot (\x,{\x*\x+2*\x+3});
\end{tikzpicture}
\end{minipage}
\end{center}

\begin{methodebox}[Méthode 6 --- Tracer une parabole à partir d'un tableau de valeurs]
\etapeexemples
Dresser un tableau de valeurs, puis tracer la courbe de :
\begin{questions}[label=\alph*)]
  \item $f(x)=-x^{2}-2x+3$ sur $\intff{-4}{2}$ ;
  \item $g(x)=x^{2}-4x$ sur $\intff{-1}{5}$.
\end{questions}
\begin{correction}
\cas{a)}\par\nopagebreak
\begin{minipage}[c]{0.58\linewidth}\raggedright
\explique{$a=-1<0$ : une « bosse ».}{le signe de $a$ servira de contrôle}
\begin{center}\small
\renewcommand{\arraystretch}{1.25}
\begin{tabular}{|c|c|c|c|c|c|c|c|}
\hline
$x$ & $-4$ & $-3$ & $-2$ & $-1$ & $0$ & $1$ & $2$\\
\hline
$f(x)$ & $-5$ & $0$ & $3$ & $4$ & $3$ & $0$ & $-5$\\
\hline
\end{tabular}
\end{center}
\explique{On place les points et on les relie à main levée, par une courbe arrondie.}{%
les valeurs se répondent de part et d'autre de $x=-1$ : symétrie}
\end{minipage}\hfill
\begin{minipage}[c]{0.38\linewidth}\centering
\begin{tikzpicture}[x=0.38cm,y=0.38cm]
  \repere{-5}{-5}{3}{5}
  \gradx{-4,-3,-2,-1,1,2} \grady{-4,-2,2,4}
  \draw[coulCourbe,line width=1.1pt,domain=-4:2,samples=60] plot (\x,{-\x*\x-2*\x+3});
  \foreach \p in {(-4,-5),(-3,0),(-2,3),(-1,4),(0,3),(1,0),(2,-5)}
    \path[coulCourbe] \p node[croix]{};
\end{tikzpicture}
\end{minipage}
\cas{b)}\par\nopagebreak
\begin{minipage}[c]{0.58\linewidth}\raggedright
\explique{$a=1>0$ : un « creux ». Racines : $x(x-4)=0$, soit $0$ et $4$.}{les points sur
l'axe des abscisses}
\begin{center}\small
\renewcommand{\arraystretch}{1.25}
\begin{tabular}{|c|c|c|c|c|c|c|c|}
\hline
$x$ & $-1$ & $0$ & $1$ & $2$ & $3$ & $4$ & $5$\\
\hline
$g(x)$ & $5$ & $0$ & $-3$ & $-4$ & $-3$ & $0$ & $5$\\
\hline
\end{tabular}
\end{center}
\end{minipage}\hfill
\begin{minipage}[c]{0.38\linewidth}\centering
\begin{tikzpicture}[x=0.38cm,y=0.38cm]
  \repere{-2}{-5}{6}{5}
  \gradx{-1,1,2,3,4,5} \grady{-4,-2,2,4}
  \draw[coulCourbe,line width=1.1pt,domain=-1:5,samples=60] plot (\x,{\x*\x-4*\x});
  \foreach \p in {(-1,5),(0,0),(1,-3),(2,-4),(3,-3),(4,0),(5,5)}
    \path[coulCourbe] \p node[croix]{};
\end{tikzpicture}
\end{minipage}
\end{correction}

\etape{À vous de jouer}
Dresser un tableau de valeurs, puis tracer la courbe de $h(x)=-x^{2}+4x$ sur
$\intff{-1}{5}$. Où coupe-t-elle l'axe des abscisses ?
\reponse{$h(-1)=-5$, $h(0)=0$, $h(1)=3$, $h(2)=4$, $h(3)=3$, $h(4)=0$, $h(5)=-5$ ; elle
coupe l'axe en $0$ et en $4$.}
\end{methodebox}

\afaire{exercices 21 à 24}

\partiecours{Sommet et variations de la parabole}

\begin{proprietebox}[Théorème 3 --- Sommet et variations]
La parabole de $f(x)=ax^{2}+bx+c$ a pour \textbf{sommet} le point
\[
  \boxed{S(\alpha\,;\beta) \quad\text{avec}\quad \alpha=-\dfrac{b}{2a}
  \quad\text{et}\quad \beta=f(\alpha)}
\]
La droite d'équation $x=\alpha$ est son axe de symétrie.
\begin{center}
\begin{minipage}{0.47\linewidth}\centering
{\small $a>0$ : $\beta$ est le \textbf{minimum}}\par\smallskip
\begin{tikzpicture}[scale=0.9]
  \tkzTabInit[lgt=2,espcl=2.4]{$x$/1,$f(x)$/2}{$-\infty$,$\alpha$,$+\infty$}
  \tkzTabVar{+/$+\infty$,-/$\beta$,+/$+\infty$}
\end{tikzpicture}
\end{minipage}\hfill
\begin{minipage}{0.47\linewidth}\centering
{\small $a<0$ : $\beta$ est le \textbf{maximum}}\par\smallskip
\begin{tikzpicture}[scale=0.9]
  \tkzTabInit[lgt=2,espcl=2.4]{$x$/1,$f(x)$/2}{$-\infty$,$\alpha$,$+\infty$}
  \tkzTabVar{-/$-\infty$,+/$\beta$,-/$-\infty$}
\end{tikzpicture}
\end{minipage}
\end{center}
\end{proprietebox}

\begin{remarquebox}
Quand le trinôme a deux racines, l'axe de symétrie passe par leur milieu :
$\alpha=\dfrac{x_{1}+x_{2}}{2}$.
\end{remarquebox}

\begin{methodebox}[Méthode 7 --- Trouver le sommet, dresser le tableau de variations]
\etapeexemples
Donner le sommet de la parabole et le tableau de variations de :
\begin{questions}[label=\alph*)]
  \item $f(x)=2x^{2}-8x+3$
  \item $g(x)=-3x^{2}-6x+4$
\end{questions}
\begin{correction}
\cas{a)}
\begin{calculs}
\alpha &= -\frac{-8}{2\times 2} &\qquad& \rem{$\alpha=-\frac{b}{2a}$, avec $a=2$ et $b=-8$}\\
\alpha &= 2 &&\\
\beta &= 2\times 2^{2}-8\times 2+3 &\qquad& \rem{$\beta=f(\alpha)$ : on remplace $x$ par $2$}\\
\beta &= -5 &&
\end{calculs}
\begin{center}
\begin{tikzpicture}[scale=0.9]
  \tkzTabInit[lgt=2,espcl=2.4]{$x$/1,$f(x)$/2}{$-\infty$,$2$,$+\infty$}
  \tkzTabVar{+/$+\infty$,-/$-5$,+/$+\infty$}
\end{tikzpicture}
\end{center}
\explique{Sommet $S(2\,;-5)$ : $f$ admet un minimum, $-5$, atteint en $x=2$.}{$a=2>0$ :
un creux, donc un minimum}
\cas{b)}
\begin{calculs}
\alpha &= -\frac{-6}{2\times(-3)} &&\\
\alpha &= -1 &\qquad& \rem{$-\frac{-6}{-6}=-1$ : attention aux signes}\\
\beta &= -3\times(-1)^{2}-6\times(-1)+4 &&\\
\beta &= 7 &&
\end{calculs}
\begin{center}
\begin{tikzpicture}[scale=0.9]
  \tkzTabInit[lgt=2,espcl=2.4]{$x$/1,$g(x)$/2}{$-\infty$,$-1$,$+\infty$}
  \tkzTabVar{-/$-\infty$,+/$7$,-/$-\infty$}
\end{tikzpicture}
\end{center}
\explique{Sommet $S(-1\,;7)$ : $g$ admet un maximum, $7$, atteint en $x=-1$.}{$a=-3<0$ :
une bosse}
\end{correction}

\etape{À vous de jouer}
Donner le sommet et le tableau de variations de :
\[
  f_{1}(x)=x^{2}+6x+4 \qquad f_{2}(x)=-x^{2}+4x+1 \qquad f_{3}(x)=2x^{2}-4x+7
\]
\reponse{$f_{1}$ : $S(-3\,;-5)$, minimum ;\; $f_{2}$ : $S(2\,;5)$, maximum ;\;
$f_{3}$ : $S(1\,;5)$, minimum.}
\end{methodebox}

\afaire{exercices 25 à 28}

\partiecours{Synthèse : étudier et tracer une parabole}

L'étude complète d'un trinôme, demandée dans les exercices et les devoirs, rassemble tout
ce qui précède : sens de la parabole, sommet, racines, variations, puis tracé.

\begin{methodebox}[Méthode 8 --- Étude complète d'une fonction du second degré]
\etapeexemples
Étudier la fonction $f(x)=-x^{2}+x+2$ sur $\intff{-2}{3}$, puis tracer sa courbe.
\begin{correction}
\cas{Résolution}\par\nopagebreak
\explique{$a=-1<0$ : une bosse.}{\textbf{sens}}
\begin{calculs}
\alpha &= -\frac{1}{2\times(-1)}=\frac{1}{2} &\qquad& \rem{\textbf{sommet}}\\
\beta &= -\frac{1}{4}+\frac{1}{2}+2=\frac{9}{4} &\qquad& \rem{$\beta=f\left(\frac{1}{2}\right)=2{,}25$}\\
\Delta &= 1^{2}-4\times(-1)\times 2=9 &\qquad& \rem{\textbf{racines}}\\
x_{1} &= \frac{-1-3}{-2}=2 &&\\
x_{2} &= \frac{-1+3}{-2}=-1 &&
\end{calculs}
\begin{minipage}[c]{0.58\linewidth}\centering
\begin{tikzpicture}[scale=0.85]
  \tkzTabInit[lgt=1.6,espcl=2]{$x$/1,$f(x)$/2}{$-2$,$\frac{1}{2}$,$3$}
  \tkzTabVar{-/$-4$,+/$\frac{9}{4}$,-/$-4$}
\end{tikzpicture}

\smallskip
{\small
\renewcommand{\arraystretch}{1.25}
\begin{tabular}{|c|c|c|c|c|c|c|}
\hline
$x$ & $-2$ & $-1$ & $0$ & $1$ & $2$ & $3$\\
\hline
$f(x)$ & $-4$ & $0$ & $2$ & $2$ & $0$ & $-4$\\
\hline
\end{tabular}}

\smallskip
\explique{On place les points et le sommet, puis on relie, symétriquement par rapport à
$x=\frac{1}{2}$.}{\textbf{tracé} (méthode 6)}
\end{minipage}\hfill
\begin{minipage}[c]{0.38\linewidth}\centering
\begin{tikzpicture}[x=0.42cm,y=0.42cm]
  \repere{-3}{-5}{4}{4}
  \gradx{-2,-1,1,2,3} \grady{-4,-2,2,3}
  \draw[black!40,dashed] (0.5,-5) -- (0.5,4);
  \draw[coulCourbe,line width=1.1pt,domain=-2:3,samples=60] plot (\x,{-\x*\x+\x+2});
  \foreach \p in {(-2,-4),(-1,0),(0,2),(1,2),(2,0),(3,-4)} \path[coulCourbe] \p node[croix]{};
  \path[coulExo] (0.5,2.25) node[croix]{} node[above,font=\scriptsize] {$S$};
\end{tikzpicture}
\end{minipage}
\end{correction}

\etape{À vous de jouer}
Faire la même étude pour $g(x)=x^{2}-x-2$ sur $\intff{-2}{3}$, puis tracer sa courbe.
\reponse{$a>0$ ; $S\left(\frac{1}{2}\,;-\frac{9}{4}\right)$, minimum ; racines $-1$ et
$2$ ; $g(-2)=4$, $g(0)=-2$, $g(1)=-2$, $g(3)=4$.}
\end{methodebox}

\afaire{exercices 29 à 31}

\partiecours{Signe d'un trinôme}

\begin{exemplebox}[Exemple --- signe de $f(x)=-3(x-2)(x+1)$]
On dresse le tableau de signes du produit ; le facteur $-3$ est toujours négatif.
\begin{center}
\begin{tikzpicture}[scale=0.9]
  \tkzTabInit[lgt=3,espcl=2.2]{$x$/1,$-3$/1,$x+1$/1,$x-2$/1,$f(x)$/1.2}%
    {$-\infty$,$-1$,$2$,$+\infty$}
  \tkzTabLine{,-,t,-,t,-,}
  \tkzTabLine{,-,z,+,t,+,}
  \tkzTabLine{,-,t,-,z,+,}
  \tkzTabLine{,-,z,+,z,-,}
\end{tikzpicture}
\end{center}
$f(x)$ est positif entre les racines, négatif à l'extérieur : le contraire du signe de
$a=-3$ entre les racines, le signe de $a$ ailleurs.
\end{exemplebox}

\begin{proprietebox}[Théorème 4 --- Signe d'un trinôme]
\begin{itemize}
  \item Si $\Delta>0$ : le trinôme a le signe de $a$ à l'\textbf{extérieur} des racines,
        le signe contraire \textbf{entre} les racines.
  \item Si $\Delta=0$ : il a le signe de $a$, et s'annule en $x_{0}$.
  \item Si $\Delta<0$ : il a \textbf{toujours} le signe de $a$.
\end{itemize}
\end{proprietebox}

\begin{center}
\begin{minipage}{0.47\linewidth}\centering
{\small $\Delta>0$ et $a>0$}\par\smallskip
\begin{tikzpicture}[scale=0.8]
  \tkzTabInit[lgt=3.2,espcl=1.5]{$x$/1,$ax^{2}+bx+c$/1.4}%
    {$-\infty$,$x_{1}$,$x_{2}$,$+\infty$}
  \tkzTabLine{,+,z,-,z,+,}
\end{tikzpicture}
\end{minipage}\hfill
\begin{minipage}{0.47\linewidth}\centering
{\small $\Delta>0$ et $a<0$}\par\smallskip
\begin{tikzpicture}[scale=0.8]
  \tkzTabInit[lgt=3.2,espcl=1.5]{$x$/1,$ax^{2}+bx+c$/1.4}%
    {$-\infty$,$x_{1}$,$x_{2}$,$+\infty$}
  \tkzTabLine{,-,z,+,z,-,}
\end{tikzpicture}
\end{minipage}
\end{center}

\begin{methodebox}[Méthode 9 --- Dresser le tableau de signes d'un trinôme]
\etapeexemples
Dresser le tableau de signes de :
\begin{questions}[label=\alph*)]
  \item $f(x)=2x^{2}+x-3$
  \item $g(x)=-2x^{2}+5x+3$
  \item $h(x)=-x^{2}+6x-9$ \; et \; $k(x)=x^{2}-x+2$
\end{questions}
\begin{correction}
\cas{a)}
\begin{calculs}
\Delta &= 1^{2}-4\times 2\times(-3)=25 &&\\
x_{1} &= \frac{-1-5}{4}=-\frac{3}{2} &&\\
x_{2} &= \frac{-1+5}{4}=1 &&
\end{calculs}
\begin{center}
\begin{tikzpicture}[scale=0.9]
  \tkzTabInit[lgt=2.4,espcl=2.2]{$x$/1,$f(x)$/1.2}{$-\infty$,$-\frac{3}{2}$,$1$,$+\infty$}
  \tkzTabLine{,+,z,-,z,+,}
\end{tikzpicture}
\end{center}
\explique{$a=2>0$ : positif à l'extérieur des racines.}{théorème 4}
\cas{b)}
\begin{calculs}
\Delta &= 5^{2}-4\times(-2)\times 3=49 &&\\
x_{1} &= \frac{-5-7}{-4}=3 &&\\
x_{2} &= \frac{-5+7}{-4}=-\frac{1}{2} &\qquad& \rem{on range les racines : $-\frac{1}{2}<3$}
\end{calculs}
\begin{center}
\begin{tikzpicture}[scale=0.9]
  \tkzTabInit[lgt=2.4,espcl=2.2]{$x$/1,$g(x)$/1.2}{$-\infty$,$-\frac{1}{2}$,$3$,$+\infty$}
  \tkzTabLine{,-,z,+,z,-,}
\end{tikzpicture}
\end{center}
\explique{$a=-2<0$ : négatif à l'extérieur, positif entre les racines.}{}
\cas{c)}
\explique{$h$ : $\Delta=36-36=0$, $x_{0}=3$ ; $a=-1<0$ : $h$ est négatif, nul en $3$.}{%
$\Delta=0$}
\explique{$k$ : $\Delta=1-8=-7<0$ ; $a=1>0$ : $k$ est toujours positif.}{$\Delta<0$}
\end{correction}

\etape{À vous de jouer}
Dresser le tableau de signes de :
\[
  f(x)=x^{2}-9 \qquad g(x)=-x^{2}+3x-2 \qquad h(x)=2x^{2}-x+3
\]
\reponse{$f$ : racines $-3$ et $3$, positif à l'extérieur ;\; $g$ : racines $1$ et $2$,
$a<0$, positif entre les racines ;\; $h$ : $\Delta=-23<0$, toujours positif.}
\end{methodebox}

\afaire{exercices 32 à 35}

\partiecours{Inéquations du second degré}

Pour résoudre une inéquation du second degré, on met tout dans le membre de gauche, on
dresse le tableau de signes du trinôme, puis on lit les solutions : crochets fermés en
une racine pour $\le$ ou $\ge$, ouverts pour $<$ ou $>$.

\begin{methodebox}[Méthode 10 --- Résoudre une inéquation du second degré]
\etapeexemples
Résoudre dans $\R$ :
\begin{questions}[label=\alph*)]
  \item $-x^{2}+5x-4>0$
  \item $x^{2}-2x-8\le 0$
  \item $2x^{2}-x\le 3$
  \item $x^{2}+2x+5>0$
\end{questions}
\begin{correction}
\cas{a)}
\explique{$\Delta=25-16=9$ ; racines $\dfrac{-5-3}{-2}=4$ et $\dfrac{-5+3}{-2}=1$.}{}
\begin{center}
\begin{tikzpicture}[scale=0.9]
  \tkzTabInit[lgt=3,espcl=2.2]{$x$/1,$-x^{2}+5x-4$/1.2}{$-\infty$,$1$,$4$,$+\infty$}
  \tkzTabLine{,-,z,+,z,-,}
\end{tikzpicture}
\end{center}
\explique{$S=\left]1\,;4\right[$.}{$a<0$ : positif entre les racines ; inégalité
stricte, racines exclues}
\cas{b)}
\explique{$\Delta=4+32=36$ ; racines $-2$ et $4$ ; $a>0$ : négatif entre les racines.}{}
\explique{$S=\intff{-2}{4}$.}{« ou nul » : racines comprises}
\cas{c)}
\begin{calculs}
2x^{2}-x &\le 3 &&\\
2x^{2}-x-3 &\le 0 &\qquad& \rem{tout dans le membre de gauche}
\end{calculs}
\explique{$\Delta=1+24=25$ ; racines $\dfrac{1-5}{4}=-1$ et $\dfrac{1+5}{4}=\dfrac{3}{2}$
; $a>0$.}{}
\explique{$S=\left[-1\,;\tfrac{3}{2}\right]$.}{négatif entre les racines}
\cas{d)}
\explique{$\Delta=4-20=-16<0$ et $a=1>0$ : le trinôme est toujours positif ;
$S=\R$.}{pas de racine : signe de $a$ partout}
\end{correction}

\etape{À vous de jouer}
Résoudre dans $\R$ :
\[
  x^{2}-x-6<0 \qquad -x^{2}-2x+8\ge 0 \qquad 2x^{2}+x+1<0 \qquad x^{2}\ge 16
\]
\reponse{$S=\left]-2\,;3\right[$ ;\; racines $-4$ et $2$ : $S=\intff{-4}{2}$ ;\;
$\Delta=-7$ et $a>0$ : $S=\varnothing$ ;\; $x^{2}-16\ge 0$ :
$S=\left]-\infty\,;-4\right]\cup\left[4\,;+\infty\right[$.}
\end{methodebox}

\afaire{exercices 36 à 39}

\partiecours{Intersection d'une parabole et d'une droite}

\begin{proprietebox}[Propriété 4 --- Points d'intersection]
Les abscisses des points communs à la parabole $\mathcal{P}$ : $y=f(x)$ et à la droite
$\mathcal{D}$ : $y=g(x)$ sont les solutions de l'équation $f(x)=g(x)$. Il y a $2$, $1$
ou $0$ point commun, selon le signe du discriminant.
\end{proprietebox}

\begin{methodebox}[Méthode 11 --- Trouver les points d'intersection]
\etapeexemples
$\mathcal{P}$ est la parabole d'équation $y=-x^{2}+2x+3$. Trouver ses points
d'intersection avec chacune des droites
$\mathcal{D}_{1}$ : $y=x+1$, \ $\mathcal{D}_{2}$ : $y=2x+3$ \ et \
$\mathcal{D}_{3}$ : $y=2x+4$.
\begin{correction}
\cas{Avec $\mathcal{D}_{1}$}
\begin{calculs}
-x^{2}+2x+3 &= x+1 &&\\
-x^{2}+x+2 &= 0 &\qquad& \rem{$f(x)=g(x)$, puis tout à gauche}
\end{calculs}
\explique{$\Delta=1+8=9$ : deux solutions, $\dfrac{-1-3}{-2}=2$ et
$\dfrac{-1+3}{-2}=-1$.}{les abscisses des points communs}
\explique{Points $A(-1\,;0)$ et $B(2\,;3)$.}{ordonnées par l'équation de la droite :
$-1+1=0$ et $2+1=3$}
\cas{Avec $\mathcal{D}_{2}$}
\begin{calculs}
-x^{2}+2x+3 &= 2x+3 &&\\
-x^{2} &= 0 &&
\end{calculs}
\explique{Une seule solution, $0$ : la droite touche la parabole en $C(0\,;3)$.}{%
$\mathcal{D}_{2}$ est tangente à $\mathcal{P}$}
\cas{Avec $\mathcal{D}_{3}$}
\begin{calculs}
-x^{2}+2x+3 &= 2x+4 &&\\
-x^{2}-1 &= 0 &&\\
x^{2} &= -1 &&
\end{calculs}
\explique{Aucune solution : $\mathcal{D}_{3}$ ne coupe pas $\mathcal{P}$.}{un carré
n'est jamais négatif}
\begin{center}
\begin{tikzpicture}[x=0.45cm,y=0.45cm]
  \repere{-5}{-5}{5}{5}
  \gradx{-4,-3,-2,-1,1,2,3,4} \grady{-4,-2,2,4}
  \draw[coulCourbe,line width=1.1pt,domain=-2:4,samples=60] plot (\x,{-\x*\x+2*\x+3});
  \draw[coulCourbeB,line width=1.1pt,domain=-5:4] plot (\x,{\x+1});
  \draw[coulCourbeC,line width=1.1pt,domain=-4:1] plot (\x,{2*\x+3});
  \draw[coulRem,line width=1.1pt,domain=-4.5:0.5] plot (\x,{2*\x+4});
  \path (-1,0) node[croix]{} (2,3) node[croix]{} (0,3) node[croix]{};
  \node[coulCourbe,font=\footnotesize] at (3.6,-1.4) {$\mathcal{P}$};
  \node[coulCourbeB,font=\footnotesize] at (4.4,4.3) {$\mathcal{D}_{1}$};
  \node[coulCourbeC,font=\footnotesize] at (1.6,4.8) {$\mathcal{D}_{2}$};
  \node[coulRem,font=\footnotesize] at (-0.4,4.8) {$\mathcal{D}_{3}$};
\end{tikzpicture}
\end{center}
\end{correction}

\etape{À vous de jouer}
Trouver les points d'intersection, s'il y en a :
\begin{questions}
  \item $\mathcal{P}$ : $y=x^{2}+2x-1$ \; et \; $\mathcal{D}$ : $y=3x+1$ ;
  \item $\mathcal{P}$ : $y=x^{2}-4x+5$ \; et \; $\mathcal{D}$ : $y=2x-4$ ;
  \item $\mathcal{P}$ : $y=x^{2}+x+2$ \; et \; $\mathcal{D}$ : $y=x+1$.
\end{questions}
\reponse{1\textdegree) $x^{2}-x-2=0$ : $(-1\,;-2)$ et $(2\,;7)$ ;\; 2\textdegree)
$x^{2}-6x+9=0$ : un seul point, $(3\,;2)$ ;\; 3\textdegree) $x^{2}+1=0$ : aucun point.}
\end{methodebox}

\afaire{exercices 40 à 43}

\partiecours{Position relative de deux courbes}

\begin{proprietebox}[Propriété 5 --- Position relative]
Pour comparer les courbes $\mathcal{C}_{f}$ et $\mathcal{C}_{g}$, on étudie le signe de
la différence $d(x)=f(x)-g(x)$ : là où $d(x)>0$, $\mathcal{C}_{f}$ est au-dessus de
$\mathcal{C}_{g}$ ; là où $d(x)<0$, elle est au-dessous ; là où $d(x)=0$, les courbes se
coupent.
\end{proprietebox}

\begin{methodebox}[Méthode 12 --- Étudier la position relative de deux courbes]
\etapeexemples
Étudier la position relative des courbes de $f$ et de $g$ :
\begin{questions}[label=\alph*)]
  \item $f(x)=-x^{2}+6x-5$ \; et \; $g(x)=x-1$ ;
  \item $f(x)=-x^{2}+4x+2$ \; et \; $g(x)=x^{2}-2x-6$.
\end{questions}
\begin{correction}
\cas{a)}
\begin{calculs}
d(x) &= -x^{2}+6x-5-(x-1) &\qquad& \rem{on soustrait \textbf{toute} l'expression de $g$}\\
d(x) &= -x^{2}+5x-4 &&
\end{calculs}
\begin{minipage}[c]{0.58\linewidth}\raggedright
\explique{$\Delta=25-16=9$ ; racines $1$ et $4$ ; $a<0$.}{}
\begin{center}
\begin{tikzpicture}[scale=0.85]
  \tkzTabInit[lgt=1.6,espcl=1.8]{$x$/1,$d(x)$/1.2}{$-\infty$,$1$,$4$,$+\infty$}
  \tkzTabLine{,-,z,+,z,-,}
\end{tikzpicture}
\end{center}
\explique{$\mathcal{C}_{f}$ est au-dessus de la droite sur $\left]1\,;4\right[$,
au-dessous ailleurs ; elles se coupent en $(1\,;0)$ et $(4\,;3)$.}{on conclut par une
phrase}
\end{minipage}\hfill
\begin{minipage}[c]{0.38\linewidth}\centering
\begin{tikzpicture}[x=0.4cm,y=0.4cm]
  \repere{-1}{-5}{6}{5}
  \gradx{1,2,3,4,5} \grady{-4,-2,2,4}
  \draw[coulCourbe,line width=1.1pt,domain=0:6,samples=60] plot (\x,{-\x*\x+6*\x-5});
  \draw[coulCourbeB,line width=1.1pt,domain=-1:6] plot (\x,{\x-1});
  \path (1,0) node[croix]{} (4,3) node[croix]{};
  \node[coulCourbe,font=\footnotesize] at (3,4.6) {$\mathcal{C}_{f}$};
  \node[coulCourbeB,font=\footnotesize] at (5.8,3.6) {$\mathcal{C}_{g}$};
\end{tikzpicture}
\end{minipage}
\cas{b)}
\begin{calculs}
d(x) &= -x^{2}+4x+2-\left(x^{2}-2x-6\right) &&\\
d(x) &= -2x^{2}+6x+8 &\qquad& \rem{le $-$ change tous les signes de $g(x)$}
\end{calculs}
\explique{$\Delta=36+64=100$ ; racines $\dfrac{-6-10}{-4}=4$ et
$\dfrac{-6+10}{-4}=-1$ ; $a<0$.}{}
\explique{$\mathcal{C}_{f}$ est au-dessus de $\mathcal{C}_{g}$ sur $\left]-1\,;4\right[$,
au-dessous ailleurs ; elles se coupent en $x=-1$ et $x=4$.}{}
\end{correction}

\etape{À vous de jouer}
Étudier la position relative des courbes de $f$ et de $g$ :
\begin{questions}
  \item $f(x)=x^{2}+x-1$ \; et \; $g(x)=3x-1$ ;
  \item $f(x)=3x^{2}+2x-5$ \; et \; $g(x)=3x^{2}-4x+1$.
\end{questions}
\reponse{1\textdegree) $d(x)=x^{2}-2x=x(x-2)$ : $\mathcal{C}_{f}$ au-dessous sur
$\left]0\,;2\right[$, au-dessus ailleurs ;\; 2\textdegree) $d(x)=6x-6$ :
$\mathcal{C}_{f}$ au-dessus pour $x>1$, au-dessous pour $x<1$.}
\end{methodebox}

\afaire{exercices 44 à 46}

\partiecours{Problèmes du second degré}

\begin{methodebox}[Méthode 13 --- Résoudre un problème du second degré]
\etapeexemples
\begin{questions}[label=\alph*)]
  \item \textbf{Un enclos.} Contre un mur, on clôt un enclos rectangulaire avec $20$~m de
        grillage, le mur formant un côté. Quelles dimensions donnent la plus grande
        aire ?
  \item \textbf{Une fusée de modélisme.} Lancée du sol, elle est à la hauteur
        $h(t)=-5t^{2}+40t$ (en mètres) au bout de $t$ secondes. Quelle hauteur maximale
        atteint-elle ? Quand retombe-t-elle ? Combien de temps reste-t-elle à plus de
        $60$~m ?
\end{questions}
\begin{correction}
\cas{a)}\par\nopagebreak
\begin{minipage}[c]{0.62\linewidth}\raggedright
\explique{Soit $x$ (en m) un côté perpendiculaire au mur ; l'autre côté mesure
$20-2x$, avec $0<x<10$.}{\textbf{inconnue}, et ses valeurs possibles}
\end{minipage}\hfill
\begin{minipage}[c]{0.34\linewidth}\centering
\begin{tikzpicture}[scale=0.55]
  \fill[black!12] (-0.3,0) rectangle (6.3,-0.35);
  \draw[very thick,black!60] (-0.3,0) -- (6.3,0);
  \draw[coulExo,line width=1.2pt] (0,0) -- (0,2) -- (6,2) -- (6,0);
  \node[font=\scriptsize,left] at (0,1) {$x$};
  \node[font=\scriptsize,above] at (3,2) {$20-2x$};
  \node[font=\scriptsize,black!60] at (3,-0.7) {mur};
\end{tikzpicture}
\end{minipage}
\begin{calculs}
A(x) &= x(20-2x)=-2x^{2}+20x &\qquad& \rem{l'\textbf{aire} en fonction de $x$}\\
\alpha &= -\frac{20}{2\times(-2)}=5 &\qquad& \rem{$a<0$ : maximum au sommet}\\
A(5) &= -50+100=50 &&
\end{calculs}
\explique{L'enclos le plus grand mesure $5$~m sur $10$~m : $50~\text{m}^{2}$.}{%
\textbf{retour à l'énoncé} : $5$ est bien entre $0$ et $10$}
\cas{b)}
\begin{calculs}
\alpha &= -\frac{40}{2\times(-5)}=4 &\qquad& \rem{$a=-5<0$ : maximum au sommet}\\
h(4) &= -80+160=80 &&
\end{calculs}
\explique{Hauteur maximale : $80$~m, au bout de $4$~s.}{}
\begin{calculs}
-5t^{2}+40t &= 0 &&\\
-5t(t-8) &= 0 &\qquad& \rem{au sol : $h(t)=0$}
\end{calculs}
\explique{Elle retombe au bout de $8$~s.}{$t=0$ est l'instant du lancement}
\begin{calculs}
-5t^{2}+40t &> 60 &&\\
-5t^{2}+40t-60 &> 0 &&\\
t^{2}-8t+12 &< 0 &\qquad& \rem{on divise par $-5<0$ : le sens change}
\end{calculs}
\explique{Racines $2$ et $6$ : $t\in\left]2\,;6\right[$, soit $4$~s au-dessus de
$60$~m.}{$a=1>0$ : négatif entre les racines}
\end{correction}

\etape{À vous de jouer}
\begin{questions}
  \item Un terrain rectangulaire a pour périmètre $30$~m et pour aire $56~\text{m}^{2}$.
        Quelles sont ses dimensions ?
  \item Un projectile est à la hauteur $h(t)=-5t^{2}+10t+15$ (en m) au bout de $t$
        secondes. Quelle est sa hauteur maximale ? Quand touche-t-il le sol ?
\end{questions}
\reponse{1\textdegree) $x(15-x)=56$, soit $x^{2}-15x+56=0$ : $8$~m sur $7$~m ;\;
2\textdegree) maximum $20$~m au bout de $1$~s ; $t^{2}-2t-3=0$ donne $3$ (on rejette
$-1$) : au sol au bout de $3$~s.}
\end{methodebox}

\afaire{exercices 47 à 49}

\end{document}
